\documentclass[12pt, a4paper]{article}

% Paquetes requeridos para el circuito
\usepackage[utf8]{inputenc}
\usepackage[T1]{fontenc}
\usepackage{circuitikz} 
\usepackage{amsmath}

\begin{document}

\begin{center}
    \begin{circuitikz}[american, scale=0.9, transform shape]
        
        % ==================== PARTE A ====================
        % 1. Fuente V1 en serie con R1
        \draw (0,0) to[V, l=$V_1$] (0,2) 
                    to[R, l=$R_1$] (0,4);
        
        % 2. Resistencia R2 en paralelo con el conjunto V1-R1
        \draw (2,0) to[R, l=$R_2$] (2,4);
        
        % Conexiones del paralelo superior e inferior
        \draw (0,4) -- (2,4);
        \draw (0,0) -- (2,0);
        
        % 3. Resistencia R3 saliendo hacia el nodo 'a'
        \draw (2,4) to[R, l=$R_3$, -o] (4,4) node[right] {a};
        
        % 4. Conexión del polo negativo hacia el nodo 'b'
        \draw (2,0) to[short, -o] (4,0) node[right] {b};
        
        \node at (1,-1) {\textbf{Parte A}};
        
        % ==================== PARTE B ====================
        % 5. Fuente de corriente en paralelo con resistencia equivalente
        \draw (6,0) to[I, l=$V_1/R_1$] (6,4);
        \draw (9,0) to[R, l=$R_1 \parallel R_2$] (9,4);
        
        % Conexiones del paralelo superior e inferior
        \draw (6,4) -- (9,4);
        \draw (6,0) -- (9,0);
        
        % 6. Resistencia R3 saliendo hacia el nodo 'a'
        \draw (9,4) to[R, l=$R_3$, -o] (11,4) node[right] {a};
        
        % 7. Conexión inferior hacia el nodo 'b'
        \draw (9,0) to[short, -o] (11,0) node[right] {b};
        
        \node at (8,-1) {\textbf{Parte B}};
        
    \end{circuitikz}
\end{center}

\end{document}